\subsection{模的外代数的定义}

定义 1. 设 A 是一个交换环，M 是一个 A-模。M 的外代数，记为 \(\bigwedge (\mathbf{M})\) 或 \(\operatorname {Alt}(\mathbf{M})\) 或 \(\bigwedge_{\mathbf{A}}(\mathbf{M})\) ，是 A 上的这样一个代数：张量代数 \(\mathsf{T}(\mathbf{M})\) 关于双边理想 \(\mathfrak{J}''\) （也记作 \(\mathfrak{J}_{\mathbf{M}}''\) ）——由元素 \(x\otimes x\) （其中 \(x\) 遍历 M）生成——的商代数。

由于理想 \(\mathfrak{J}''\) 由次数为 2 的齐次元素生成，它是一个分次理想（II，§11，第3号，命题2）；我们记 \(\mathfrak{J}_n'' = \mathfrak{J}'' \cap \mathsf{T}^n(\mathbf{M})\) ；代数 \(\bigwedge(\mathbf{M})\) 因此由诸 \(\bigwedge^n(\mathbf{M}) = \mathsf{T}^n(\mathbf{M}) / \mathfrak{J}_n''\) 组成的（称为典范的）分次进行分次。则 \(\mathfrak{J}_0'' = \mathfrak{J}_1'' = \{0\}\) ，因此 \(\bigwedge^0(\mathbf{M})\) 与 A 等同，且 \(\bigwedge^1(\mathbf{M})\) 与 \(\mathsf{T}^1(\mathbf{M}) = \mathbf{M}\) 等同；在下文中我们始终作这些等同，且典范单射 \(\mathbf{M} \to \bigwedge(\mathbf{M})\) 将记作 \(\phi''\) 或 \(\phi_{\mathbf{M}}''\) 。

命题 1. 设 E 为一个 A-代数，且 \(f: \mathbf{M} \to \mathbf{E}\) 为一个满足如下条件的 A-线性映射：

\[
(f (x)) ^ {2} = 0 \quad \text { for all } x \in \mathbf {M}.\tag{1}
\]

存在唯一的 A-代数同态 \(g: \bigwedge(\mathbf{M}) \to \mathbf{E}\) ，使得 \(f = g \circ \phi''\) 。

换言之， \((\bigwedge (\mathbf{M}),\phi '')\) 是泛映射问题（集合论，IV，§3，第1号）的一个解，其中 \(\Sigma\) 是 A-代数结构的类， \(\alpha\) -映射为从 A-模 M 到满足 (1) 的 A-代数的线性映射。

\(g\) 的唯一性由以下事实推出： \(\phi''(\mathbf{M}) = \mathbf{M}\) 生成 \(\bigwedge (\mathbf{M})\) 。为证明 \(g\) 的存在性，我们注意到由 § 5，第 1 号，命题 1，存在一个 A-代数同态 \(g_1: \mathsf{T}(\mathbf{M}) \to \mathbf{E}\) ，使得 \(f = g_1 \circ \phi\) ；我们需要证明 \(g_1\) 在理想 \(\mathfrak{J}''\) 上为零，因为若

\[
p \colon \mathbf {T} (\mathbf {M}) \rightarrow \bigwedge (\mathbf {M}) = \mathbf {T} (\mathbf {M}) / \mathfrak {J} ^ {\prime \prime}
\]

是典范同态，则可写 \(g_1 = g \circ p\) ，其中 \(g: \bigwedge (\mathbf{M}) \to \mathbf{E}\) 是一个代数同态，且结论将由以下事实推出： \(p \circ \phi = \phi''\) 。现在， \(g_1\) 的核是一个双边理想，由 (1) 及关系 \(g_1 \circ \phi = f\) ，它包含元素 \(x \otimes x\) （对 \(x \in \mathbf{M}\) ）。这就完成了证明。

注记. (1) 设 E 为 Z 型分次 A-代数，其分次为 \((\mathbf{E}_{n})\) ，并设线性映射 f（假定满足 (1)）满足

\[
f (\mathbf {M}) \subset \mathbf {E} _ {1}.\tag{2}
\]

则关系 \(g(x_1x_2 \ldots x_p) = f(x_1)f(x_2) \ldots f(x_p)\) （其中 \(x_i \in \mathbf{M}\) ）表明 \(g(\bigwedge^p (\mathbf{M})) \subset \mathrm{E}_p\) （对所有 \(p \geqslant 0\) ），因此 \(g\) 是一个分次代数同态。

\begin{enumerate}
\def\labelenumi{(\arabic{enumi})}
\setcounter{enumi}{1}
\tightlist
\item
  为避免混淆，外代数 \(\bigwedge(\mathbf{M})\) 中两个元素 u, v 的乘积通常记为 \(u \wedge v\) ，并称为 u 与 v 的外积。 \(\bigwedge^{n}(\mathbf{M})\) 的元素因此是形如 \(x_{1} \wedge x_{2} \wedge \cdots \wedge x_{n}\) 的元素之和（其中 \(x_{i} \in M\) ， \(1 \leqslant i \leqslant n\) ），常称为 n-向量。
\end{enumerate}

